\documentclass[11pt]{article}
% Shared preamble for the PNPInverse onboarding docs.
% Each document does % Shared preamble for the PNPInverse onboarding docs.
% Each document does % Shared preamble for the PNPInverse onboarding docs.
% Each document does \input{preamble} after \documentclass.

\usepackage[utf8]{inputenc}
\usepackage[T1]{fontenc}
\usepackage{lmodern}
\usepackage[margin=1in]{geometry}
\usepackage{amsmath, amssymb, amsthm}
\usepackage{mathtools}
\usepackage{empheq}
\usepackage{siunitx}
\usepackage{booktabs}
\usepackage{array}
\usepackage{enumitem}
\usepackage{xcolor}
\usepackage{framed}
\usepackage{microtype}
\usepackage{listings}
\usepackage{tikz}
\usetikzlibrary{arrows.meta, positioning, decorations.pathreplacing, calligraphy}
\usepackage[hidelinks]{hyperref}

% --- branding -------------------------------------------------------------
\definecolor{pnpblue}{RGB}{30,70,120}
\definecolor{pnpgray}{RGB}{90,90,90}
\definecolor{calloutbg}{RGB}{245,245,235}
\definecolor{shadecolor}{RGB}{245,245,235}

% --- callouts (uses framed.sty's shaded environment) ----------------------
\newenvironment{callout}[1]
  {\begin{shaded}\noindent\textbf{\color{pnpblue}#1}\par\smallskip\ignorespaces}
  {\end{shaded}}

\newcommand{\concept}[1]{\textbf{\color{pnpblue}#1}}
\newcommand{\warn}[1]{\textbf{\color{red!70!black}#1}}

% --- code listings --------------------------------------------------------
\lstset{
  basicstyle=\ttfamily\small,
  breaklines=true,
  columns=fullflexible,
  keepspaces=true,
  showstringspaces=false,
  frame=single,
  framesep=4pt,
  backgroundcolor=\color{calloutbg},
  rulecolor=\color{pnpgray!40},
}

% --- math shortcuts -------------------------------------------------------
\newcommand{\R}{\mathbb{R}}
\newcommand{\diff}{\mathrm{d}}
\newcommand{\eqd}{\overset{!}{=}}
\newcommand{\grad}{\nabla}
\newcommand{\divg}{\nabla\!\cdot}
\newcommand{\eqdef}{\;\coloneqq\;}
\newcommand{\evalat}[2]{\left.#1\right\rvert_{#2}}
\newcommand{\set}[1]{\{#1\}}

% common chemistry / physics symbols
\newcommand{\Hyd}{\mathrm{H}}
\newcommand{\Oxy}{\mathrm{O}}
\newcommand{\HtwoO}{\mathrm{H_2O}}
\newcommand{\HtwoOtwo}{\mathrm{H_2O_2}}
\newcommand{\Otwo}{\mathrm{O_2}}
\newcommand{\Hp}{\mathrm{H^+}}
\newcommand{\OHm}{\mathrm{OH^-}}
\newcommand{\Kp}{\mathrm{K^+}}
\newcommand{\Csp}{\mathrm{Cs^+}}
\newcommand{\SOfour}{\mathrm{SO_4^{2-}}}
\newcommand{\Vrhe}{V_{\mathrm{RHE}}}
\newcommand{\eo}{E^{\circ}}
\newcommand{\Eeq}{E_{\mathrm{eq}}}
\newcommand{\kzero}{k_0}

\newcommand{\onboardingfooter}{%
  \vfill\noindent\rule{\linewidth}{0.4pt}\par
  {\small\color{pnpgray} PNPInverse intern onboarding,
  compiled \today. Source under \texttt{docs/onboarding/}.}\par}
 after \documentclass.

\usepackage[utf8]{inputenc}
\usepackage[T1]{fontenc}
\usepackage{lmodern}
\usepackage[margin=1in]{geometry}
\usepackage{amsmath, amssymb, amsthm}
\usepackage{mathtools}
\usepackage{empheq}
\usepackage{siunitx}
\usepackage{booktabs}
\usepackage{array}
\usepackage{enumitem}
\usepackage{xcolor}
\usepackage{framed}
\usepackage{microtype}
\usepackage{listings}
\usepackage{tikz}
\usetikzlibrary{arrows.meta, positioning, decorations.pathreplacing, calligraphy}
\usepackage[hidelinks]{hyperref}

% --- branding -------------------------------------------------------------
\definecolor{pnpblue}{RGB}{30,70,120}
\definecolor{pnpgray}{RGB}{90,90,90}
\definecolor{calloutbg}{RGB}{245,245,235}
\definecolor{shadecolor}{RGB}{245,245,235}

% --- callouts (uses framed.sty's shaded environment) ----------------------
\newenvironment{callout}[1]
  {\begin{shaded}\noindent\textbf{\color{pnpblue}#1}\par\smallskip\ignorespaces}
  {\end{shaded}}

\newcommand{\concept}[1]{\textbf{\color{pnpblue}#1}}
\newcommand{\warn}[1]{\textbf{\color{red!70!black}#1}}

% --- code listings --------------------------------------------------------
\lstset{
  basicstyle=\ttfamily\small,
  breaklines=true,
  columns=fullflexible,
  keepspaces=true,
  showstringspaces=false,
  frame=single,
  framesep=4pt,
  backgroundcolor=\color{calloutbg},
  rulecolor=\color{pnpgray!40},
}

% --- math shortcuts -------------------------------------------------------
\newcommand{\R}{\mathbb{R}}
\newcommand{\diff}{\mathrm{d}}
\newcommand{\eqd}{\overset{!}{=}}
\newcommand{\grad}{\nabla}
\newcommand{\divg}{\nabla\!\cdot}
\newcommand{\eqdef}{\;\coloneqq\;}
\newcommand{\evalat}[2]{\left.#1\right\rvert_{#2}}
\newcommand{\set}[1]{\{#1\}}

% common chemistry / physics symbols
\newcommand{\Hyd}{\mathrm{H}}
\newcommand{\Oxy}{\mathrm{O}}
\newcommand{\HtwoO}{\mathrm{H_2O}}
\newcommand{\HtwoOtwo}{\mathrm{H_2O_2}}
\newcommand{\Otwo}{\mathrm{O_2}}
\newcommand{\Hp}{\mathrm{H^+}}
\newcommand{\OHm}{\mathrm{OH^-}}
\newcommand{\Kp}{\mathrm{K^+}}
\newcommand{\Csp}{\mathrm{Cs^+}}
\newcommand{\SOfour}{\mathrm{SO_4^{2-}}}
\newcommand{\Vrhe}{V_{\mathrm{RHE}}}
\newcommand{\eo}{E^{\circ}}
\newcommand{\Eeq}{E_{\mathrm{eq}}}
\newcommand{\kzero}{k_0}

\newcommand{\onboardingfooter}{%
  \vfill\noindent\rule{\linewidth}{0.4pt}\par
  {\small\color{pnpgray} PNPInverse intern onboarding,
  compiled \today. Source under \texttt{docs/onboarding/}.}\par}
 after \documentclass.

\usepackage[utf8]{inputenc}
\usepackage[T1]{fontenc}
\usepackage{lmodern}
\usepackage[margin=1in]{geometry}
\usepackage{amsmath, amssymb, amsthm}
\usepackage{mathtools}
\usepackage{empheq}
\usepackage{siunitx}
\usepackage{booktabs}
\usepackage{array}
\usepackage{enumitem}
\usepackage{xcolor}
\usepackage{framed}
\usepackage{microtype}
\usepackage{listings}
\usepackage{tikz}
\usetikzlibrary{arrows.meta, positioning, decorations.pathreplacing, calligraphy}
\usepackage[hidelinks]{hyperref}

% --- branding -------------------------------------------------------------
\definecolor{pnpblue}{RGB}{30,70,120}
\definecolor{pnpgray}{RGB}{90,90,90}
\definecolor{calloutbg}{RGB}{245,245,235}
\definecolor{shadecolor}{RGB}{245,245,235}

% --- callouts (uses framed.sty's shaded environment) ----------------------
\newenvironment{callout}[1]
  {\begin{shaded}\noindent\textbf{\color{pnpblue}#1}\par\smallskip\ignorespaces}
  {\end{shaded}}

\newcommand{\concept}[1]{\textbf{\color{pnpblue}#1}}
\newcommand{\warn}[1]{\textbf{\color{red!70!black}#1}}

% --- code listings --------------------------------------------------------
\lstset{
  basicstyle=\ttfamily\small,
  breaklines=true,
  columns=fullflexible,
  keepspaces=true,
  showstringspaces=false,
  frame=single,
  framesep=4pt,
  backgroundcolor=\color{calloutbg},
  rulecolor=\color{pnpgray!40},
}

% --- math shortcuts -------------------------------------------------------
\newcommand{\R}{\mathbb{R}}
\newcommand{\diff}{\mathrm{d}}
\newcommand{\eqd}{\overset{!}{=}}
\newcommand{\grad}{\nabla}
\newcommand{\divg}{\nabla\!\cdot}
\newcommand{\eqdef}{\;\coloneqq\;}
\newcommand{\evalat}[2]{\left.#1\right\rvert_{#2}}
\newcommand{\set}[1]{\{#1\}}

% common chemistry / physics symbols
\newcommand{\Hyd}{\mathrm{H}}
\newcommand{\Oxy}{\mathrm{O}}
\newcommand{\HtwoO}{\mathrm{H_2O}}
\newcommand{\HtwoOtwo}{\mathrm{H_2O_2}}
\newcommand{\Otwo}{\mathrm{O_2}}
\newcommand{\Hp}{\mathrm{H^+}}
\newcommand{\OHm}{\mathrm{OH^-}}
\newcommand{\Kp}{\mathrm{K^+}}
\newcommand{\Csp}{\mathrm{Cs^+}}
\newcommand{\SOfour}{\mathrm{SO_4^{2-}}}
\newcommand{\Vrhe}{V_{\mathrm{RHE}}}
\newcommand{\eo}{E^{\circ}}
\newcommand{\Eeq}{E_{\mathrm{eq}}}
\newcommand{\kzero}{k_0}

\newcommand{\onboardingfooter}{%
  \vfill\noindent\rule{\linewidth}{0.4pt}\par
  {\small\color{pnpgray} PNPInverse intern onboarding,
  compiled \today. Source under \texttt{docs/onboarding/}.}\par}


\title{\bfseries PNPInverse \\
       \large Project Overview \& Motivation}
\author{Onboarding doc 00 of 06}
\date{\today}

\begin{document}
\maketitle

\begin{callout}{What this document is for}
A bird's-eye view of the project: the scientific question we are
trying to answer, what kind of code we write to answer it, where you
fit in, and what the rest of the onboarding documents will cover.
You should be able to read this in 15--20 minutes. The next docs go
deep into the chemistry, the math, and the code.
\end{callout}

\tableofcontents
\bigskip\hrule\bigskip

% =========================================================================
\section{The big picture in one paragraph}

We are building a computational model of an
\concept{electrochemical cell} --- specifically the surface where an
electrode meets a salt-water electrolyte and the
\concept{oxygen reduction reaction (ORR)} happens. The ORR is how
oxygen gas, $\Otwo$, gets reduced to either water ($\HtwoO$) or
hydrogen peroxide ($\HtwoOtwo$). Which product you get, and at what
voltage, depends on a tangled mix of catalyst surface chemistry,
local pH, and how ions arrange themselves in the few nanometers
nearest the electrode. Our group at Northwestern (the Mangan and
Seitz labs) has experimental data showing that the choice of
\concept{spectator cation} --- the ``inert'' positive ion in the
salt, like $\Kp$, $\Csp$, $\mathrm{Na^+}$, $\mathrm{Li^+}$ ---
strongly changes which ORR product the cell makes, even though the
cation itself doesn't appear in the chemical equation. We want to
build a forward simulator that reproduces this cation-dependent
behavior from physical principles, and then (later) run that
simulator backward to extract the underlying kinetic constants.

\section{Why this matters scientifically}

\begin{itemize}[leftmargin=*]
\item \concept{$\HtwoOtwo$ is a useful product.} Hydrogen peroxide is
  a green oxidant used in water treatment, paper bleaching, and
  potentially as a fuel. Today industrial $\HtwoOtwo$ is made by the
  energy-intensive anthraquinone process. Electrochemical $\HtwoOtwo$
  generation from $\Otwo$ + electricity is a candidate replacement,
  but only if you can hit it \textit{selectively} --- if your
  electrode reduces $\Otwo$ all the way to water you get nothing
  useful.
\item \concept{Cation effects are mysterious and reproducible.} For
  decades, electrochemists have noticed that swapping
  $\mathrm{Li^+} \leftrightarrow \mathrm{Cs^+}$ in an ORR cell can
  shift current densities and selectivities by factors of 2--10.
  Modern theories (Co \& Zhang 2019, Singh 2016, Ringe 2019) blame
  this on cations modulating the electric double layer and the local
  pH near the electrode, not on cations entering the reaction. There
  is no first-principles forward model that nails this quantitatively.
\item \concept{Inverse parameter inference is the long game.} The
  rate constants for ORR ($\kzero$, transfer coefficient $\alpha$, etc.)
  cannot be measured directly --- they are inferred by fitting an
  assumed kinetic model to current-voltage curves. Most groups fit
  toy models (no spatial transport, no double layer). Ours is one of
  very few attempts to fit a full PDE-based model with adjoint
  gradients, which is what gives the inverse problem any hope of being
  identifiable. \warn{This pipeline is currently paused} until the
  forward model is mature enough to be trusted; that's what we are
  working on now.
\end{itemize}

\section{What kind of code we write}

\subsection{The \emph{forward} problem}

Given:
\begin{itemize}[leftmargin=*]
  \item a 1D or 2D slice of the electrode--electrolyte region
        (typically $\sim$10--100 $\mu$m thick),
  \item bulk concentrations of $\Otwo$, $\Hp$, $\OHm$, $\Kp$,
        $\SOfour$,
  \item a set of \concept{kinetic parameters}
        $\{\kzero^{(2e)}, \kzero^{(4e)}, \alpha^{(2e)}, \alpha^{(4e)}\}$
        for the two parallel ORR reactions,
  \item a value of the applied electrode potential $\Vrhe$,
\end{itemize}
the forward solver computes the steady-state distribution of every
species concentration and the electric potential everywhere in the
domain, and from those it computes the \concept{disk current density}
(total current at the electrode) and the \concept{peroxide current
density} (the slice of that current going into $\HtwoOtwo$).

The math underneath is the \concept{Poisson--Nernst--Planck (PNP)}
system coupled to a \concept{Butler--Volmer (BV)} boundary condition
at the electrode. Doc 02 derives both. Doc 03 explains how we
discretize and solve them.

\subsection{The \emph{inverse} problem (paused)}

Given experimental current-voltage data from a real electrochemistry
experiment, find the kinetic parameters that, when fed into the
forward solver, reproduce the data. This is a nonlinear PDE-constrained
optimization problem; we solve it with adjoint gradients (each forward
solve is expensive, so we cannot afford finite-difference Jacobians).
The math is interesting but the pipeline is paused until the forward
model is trustworthy enough that fitting it is meaningful. You will
not touch this in your first few months.

\section{The experimental anchor: Seitz / Mangan / Ruggiero 2022}

Everything we model is grounded in real data from our collaborators.
The canonical paper is:

\begin{quote}
\small
Ruggiero, Sanroman Gutierrez, George, Mangan, Notestein, Seitz.
\textit{Probing the Relationship Between Bulk and Local Environments
to Understand Impacts on Electrocatalytic Oxygen Reduction Reaction.}
\textbf{Journal of Catalysis}, 2022.
\end{quote}

Niall Mangan (one of the co-authors) is the PI of our lab on the
applied-math side. Linsey Seitz's group does the experiments. The
paper measures ORR currents on a CMK-3 mesoporous-carbon electrode
in $\mathrm{K_2SO_4}$ electrolyte at bulk pH 4--6, using a
\concept{rotating ring--disk electrode (RRDE)}. The disk current
tells you how fast oxygen is being consumed; the ring current is
calibrated to tell you how much $\HtwoOtwo$ leaves the disk and gets
re-oxidized at the (downstream) ring. Together they let you extract
selectivity (\% peroxide) and effective electron number $n_e$.

The 2022 paper is the peer-reviewed source. A 2025 conference deck
extends the dataset to four cations ($\mathrm{Li^+}$,
$\mathrm{Na^+}$, $\mathrm{K^+}$, $\mathrm{Cs^+}$) at the same
electrochemistry. The deck is the moving target our forward model is
chasing.

You should at some point read the Ruggiero 2022 PDF
(\texttt{docs/papers/Ruggiero2022\_JCatal\_manuscript.pdf}) and skim
the doc \texttt{docs/papers/Ruggiero2022\_JCatal\_source\_paper.md}
which is our internal cheat-sheet of the paper's hard numbers (the
collection efficiency $N=0.224$, the rotation rate 1600 rpm, the
electrolyte composition, etc.). Doc 01 in this onboarding series is
basically a chemistry tour of the same physics.

\section{Where the project is right now (May 2026)}

The codebase has gone through several major versions. The current
production stack is sometimes called the
\concept{three-species log-c muh stack} (don't worry about the
acronyms yet --- doc 03 unpacks them). It models:

\begin{itemize}[leftmargin=*]
  \item three dynamic species: $\Otwo$, $\HtwoOtwo$, $\Hp$;
  \item one or two analytic counter-ions ($\Kp$, $\SOfour$) computed
        from a closed-form Boltzmann-with-steric-saturation expression;
  \item parallel 2-electron and 4-electron ORR pathways at their
        Ruggiero 2022 standard potentials ($\eo = 0.695$ V and 1.23 V
        vs.\ RHE respectively);
  \item a finite \concept{Stern compact-layer} capacitance at the
        electrode surface ($C_S = 0.20$ F/m$^2$);
  \item optionally, water self-ionization ($\HtwoO \rightleftharpoons
        \Hp + \OHm$) and \concept{cation hydrolysis at the OHP}
        (Singh 2016 field-dependent pK$_a$).
\end{itemize}

We are deep in \concept{Phase 6$\beta$}, which is the project name
for adding cation hydrolysis. The most recent landed work, as of
2026-05-11, is ``Step 9 B.2'' --- a 14-rung sweep over hydrolysis
rate $\times$ activation parameter at a single calibration voltage.
Doc 04 chronicles the Phase~6 history; you don't need to memorize
it but the language (``Phase 6$\alpha$'', ``v10b'', ``Step 9'',
``Phase D'') comes up constantly in commits and in the Slack-style
handoff notes under \texttt{docs/handoffs/}.


\section{Map of the rest of these docs}

\begin{description}[leftmargin=*]
  \item[Doc 01 --- Electrochemistry primer.] Everything you need
    about electrodes, the electric double layer, Butler--Volmer
    kinetics, and the ORR reaction set. This is the longest doc.
  \item[Doc 02 --- The PNP--BV mathematical model.] How the PDEs are
    written and what each term means. Limited PDE background
    assumed.
  \item[Doc 03 --- Numerics and codebase tour.] Newton, continuation,
    Firedrake, and a guided tour of the most-edited files.
  \item[Doc 04 --- Phase 6$\beta$ status and open questions.]
    Where the science is right now, what's blocked, and what kinds
    of contributions are useful for a new contributor.
  \item[Doc 05 --- Meeting agenda \& getting started.]
    The plan for our first meeting and a hands-on setup checklist.
\end{description}

\onboardingfooter
\end{document}
